\chapter*{第二部分导读：把数据类型读成科学证据}
\addcontentsline{toc}{chapter}{第二部分导读：把数据类型读成科学证据}

Part II 是本书从“科研计算基础”进入“科学数据处理”的转折点。从这里开始，我们不再只问“代码能不能跑”，而要反复追问：这份数据是什么物理量的记录？它的单位、误差、坐标、采样和选择条件是什么？经过处理之后，它还能支持怎样的科学判断？

天文和物理数据的外观差异很大。星表像普通表格，图像像二维矩阵，光谱是一维曲线，时间序列带有采样窗口，实验和模拟数据又常常被模型方程约束。但在教材层面，它们其实共享一条主线：

\begin{enumerate}
    \item 先确认数据对象和字段契约；
    \item 再确认单位、误差和测量条件；
    \item 根据数据类型选择合适的最小处理算法；
    \item 用图表和残差检查处理结果是否合理；
    \item 最后把结果写成有边界的科学解释。
\end{enumerate}

如果说 Part I 训练的是“把项目做得可运行、可记录、可复查”，那么 Part II 训练的就是“把数据读成带物理含义的证据”。

\section*{本部分为什么先讲误差和单位}

Part II 的第一章不是星表、图像或光谱，而是科学数据、误差与单位。这不是顺序上的偶然。因为所有数据类型最终都要回到同一个问题：你得到的数值是否能被解释，解释的可信度有多高，在哪些条件下会失效。

例如：

\begin{itemize}
    \item 视差反演距离时，低信噪比会放大不确定度；
    \item 孔径测光时，背景估计会直接影响净信号；
    \item 谱线测量时，归一化和采样会改变等效宽度；
    \item 周期搜索时，采样窗口可能制造假周期；
    \item 实验拟合时，残差结构会暴露模型假设的失败。
\end{itemize}

这些问题表面上属于不同章节，背后却都在训练同一种判断：不要把“算出来的值”直接当成“可靠的科学事实”。

\section*{五类数据的共同阅读框架}

为了避免每章学完之后各自散掉，本部分建议你带着同一组问题阅读每种数据：

\begin{enumerate}
    \item \textbf{对象是什么}：一行、一个像素、一条谱线、一个时间点或一次实验记录分别代表什么；
    \item \textbf{坐标是什么}：横轴、纵轴、像素坐标、天球坐标、波长轴或时间轴如何定义；
    \item \textbf{单位和误差是什么}：数值能否相加、比较、拟合，误差是否会传播；
    \item \textbf{处理算法做了什么}：筛选、派生、背景扣除、归一化、折叠或拟合改变了什么；
    \item \textbf{解释边界在哪里}：当前结果能支持什么，不能支持什么，哪些选择效应还没有处理。
\end{enumerate}

\begin{protip}[统一问题，不统一答案]
Part II 的目标不是把所有数据都强行化成同一种格式，而是训练一种统一的提问方式。数据类型可以不同，但你对单位、误差、采样、元数据和解释边界的敏感性应当越来越稳定。
\end{protip}

\section*{本部分的结果报告底线}

从本部分开始，任何一个“算出来的结果”都不应孤立出现。一个 HR 图、一组孔径测光结果、一条红移估计、一个候选周期或一次模型拟合，至少都应回答五个问题：用了哪些数据，采用了哪条公式或算法，结果的单位和误差怎样理解，哪些图表或残差支持它，以及它在哪些条件下不能被过度解释。

这条底线会反复出现在后面的章节里。它的目的不是增加写作负担，而是帮助你养成一种科研表达习惯：\textbf{结果必须带着来源、方法、诊断和限制一起交付。} 只有这样，数值才不会从科学证据退化成没有上下文的输出。

\section*{Part II 如何接向机器学习}

后面的机器学习章节会把数据变成特征、标签、训练集、验证集和模型输入。Part II 的作用，就是在此之前让你知道这些输入从哪里来。

如果没有 Part II 的训练，机器学习很容易退化成“把表丢给模型”。而科学项目真正需要的是：

\begin{itemize}
    \item 知道某个特征是不是可靠派生量；
    \item 知道某个标签是否来自清楚的物理或人工判读过程；
    \item 知道样本筛选是否已经改变了总体分布；
    \item 知道图像、光谱或时间序列在进入模型前经历了哪些归一化和重采样；
    \item 知道模型错误是否可能反映数据质量问题，而不仅是算法问题。
\end{itemize}

因此，Part II 不是机器学习前的“准备工作”，而是机器学习可信度的一部分。
